

 






\section{Evaluation}
\label{sec:exps}

We evaluate \ours on \numtasks diverse tasks: 
Dialogue Response Generation \citep[\Cref{section:dialogue}; ][]{mehri-eskenazi-2020-unsupervised}, 
Code Optimization \citep[\Cref{sec:pie}; ][]{pie}, 
Code Readability Improvement \citep[\Cref{section:code}; ][]{puri2021codenet}, 
Math Reasoning \citep[\Cref{section:gsm}; ][]{cobbe2021training}, 
Sentiment Reversal \citep[\Cref{section:sentiment}; ][]{zhang2015character}, 
and we introduce two new tasks: Acronym Generation (\Cref{section:acronym}) and 
\commongenhard (a harder version of \citet{lin2019commongen} with 20-30 keyword constraints instead of 3-5; \Cref{section:constrainedgen})

Examples for all tasks and dataset statistics are provided in \autoref{tab:task_descriptions} (\Cref{app:tasks}). 

\subsection{Instantiating \ours} %
\label{ssec:instantiation}
We instantiate \ours following the high-level description in \Cref{sec:method}. %
The \fbmod-\itermod iterations continue until the desired output quality or task-specific criterion is reached, 
up to a maximum of 4 iterations.
To make our evaluation consistent across different models, we implemented both \fbmod and \itermod as few-shot prompts even with models that respond well to instructions, such as \chatgpt and \gptf. 


\paragraph{Base LLMs} 
Our main goal is to evaluate whether we can improve the performance of any strong base LLMs using \ours. 
Therefore, we compare \ours to the same base LLMs but without feedback-refine iterations.
We used three main strong base LLM across all tasks: 
\gptt (\texttt{text-davinci-003}), \chatgpt (\texttt{gpt-3.5-turbo}), and \gptf \citep{openai2023gpt4}.
For code-based tasks, we also experimented with \codex (\texttt{code-davinci-002}).
In all tasks, either \gptt or \gptf is the previous state-of-the-art.\footnote{A comparison with other few-shot and fine-tuned approaches is provided in \Cref{sec:sota}}
We used the same prompts from previous work when available (such as for \codeopt and \gsm); otherwise, we created  prompts as detailed in \Cref{sec:prompts}.
We use greedy decoding with a temperature of 0.7 for all setups.    












